\SourceLeaf{184}\clearpage
\begin{center}\textsc{Gauss} an \textsc{Olbers}. Göttingen, 28. April 1817.\end{center}

.\ .\ .\ .\ .\ .\ .\ \textsc{Wachter} hat eine kleine Piece drucken lassen über die
ersten Gründe der Geometrie, wovon Sie durch \textsc{Lindenau} vermuthlich
ein Exemplar erhalten werden. Obgleich W. in das Wesen der Sache mehr
eingedrungen ist, als seine Vorgänger, so ist sein Beweis doch nicht
bündiger als alle andern. Ich komme immer mehr zu der Überzeugung, dass die
Nothwendigkeit unserer Geometrie nicht bewiesen werden kann, wenigstens
nicht \textls{vom} \textbf{\textls{menschlichen}} Verstande noch \textls{für} den menschlichen Verstand.
Vielleicht kommen wir in einem andern Leben zu andern Einsichten in das
Wesen des Raums, die uns jetzt unerreichbar sind. Bis dahin müsste man die
Geometrie nicht mit der Arithmetik, die rein a priori steht, sondern etwa
mit der Mechanik in gleichen Rang setzen. .\ .\ .\ .\ .\ .

\begin{center}\rule{3cm}{0.4pt}\end{center}

